The sails came out of the south at dawn, and Euron knew them for what they were before his lookout had finished calling the count down from the masthead, because there was only one fleet left afloat in all of Westeros that would come up the Blackwater under longship sail rather than a trader’s canvas, flying no banner it needed a spyglass to read.

Twenty-odd ships. His niece’s, then. He had wondered, in an idle way, across the long months of this siege, whether Yara would come — had rather assumed she would not, having left her chained in a hold once and freed by circumstance rather than mercy, and having credited her, ever since, with enough sense to hold her rock and let the mainland’s wars burn themselves out without her. He found, watching the sails resolve out of the grey morning, that he had misjudged her, which did not trouble him nearly so much as the fact of having misjudged her at all. Euron Greyjoy did not, as a rule, misjudge the people he had already broken once. It suggested the breaking had been less thorough than he’d credited, which was a thing worth knowing about a woman sailing toward him with twenty ships and, he did not doubt, a grievance she’d spent a year and more sharpening to a point.

“Theon’s dead,” he said aloud, to no one, because his crew could not answer him and he found, sometimes, that saying a thing aloud clarified it better than turning it over in silence. Word of Winterfell had reached the city a fortnight past. If the boy had died there — and the rumors, garbled as they came, seemed sure enough of a Greyjoy dying at that godswood to be worth believing — then Yara sailed south carrying the last of the debt she reckoned Euron owed their whole drowned line, and meant, plainly, to collect it in the only coin the Iron Islands had ever properly understood.

He found the thought did not frighten him. Very little frightened him anymore, which he had come, across a long life, to regard as either his greatest strength or the surest sign of whatever had gone wrong in him somewhere too early to remember. But it clarified things, the way an enemy’s arrival always clarified things, stripping away the slow attrition of the siege — the deserters, the grumbling, Roland Storm’s careful watching face, Cersei’s endless shut gates — and replacing all of it with a single, simple question he had always found considerably easier to answer than the complicated ones.

What was he willing to spend, to still be standing when this was finished.

\sceneBreak

The Golden Company would not hold. He had known it for days, whatever he’d told Cersei about contracts and living memory, and Yara’s sails only confirmed the timing of a thing he’d already reckoned inevitable. Ten thousand sellswords who had signed to fight the living did not hold their line against an army of the dead five days out, not for a queen who hanged her own deserters in public squares, and not — this was the part Cersei had never understood, for all her cleverness — because they were cowards, but because they were professionals, and a professional understood better than any zealot exactly when a contract had stopped being worth the paper it was inked on.

Roland Storm understood it too. Euron had watched the man understand it, these past weeks, the way he watched most things — patiently, from a small calculating distance, filing away each fraction of hesitation for whenever it became useful. The captain had not yet decided which way he meant to break. But he would break; they all would; and Euron had spent a great deal of this siege quietly ensuring that when the Golden Company finally came apart, at least some meaningful fraction of it came apart in a direction that left Euron Greyjoy holding more than he’d started with rather than less.

None of which addressed the dragons.

That was the piece that had never resolved itself, across all his patient arithmetic. Cersei behind her walls, the dead at the gate, the Golden Company fracturing, Yara sailing up his own harbor with a year’s grievance — all of it he could work with, spend against, turn to his own account with time and patience and the particular ruthlessness he’d never yet found a situation that exhausted. But the dragon queen had a dragon still, and the King in the North was, if the whispers out of that broken column held any truth, blood enough to ride him too should the need ever come. And a man could not out-patience a dragon. A man could not spend a dragon against itself, or wait for a dragon to break the way a sellsword broke. There was only one thing in all the world that Euron had ever heard tell could bend a dragon to a will not its rider’s, and it had been sitting wrapped in oiled leather behind his own cabin door for the better part of two years, costing him three days of a ringing skull merely for the testing of it.

He had told himself, across all those months, that he was waiting for the right moment.

He understood now, watching his niece’s sails come up the harbor, that he had in truth been waiting to be frightened enough to stop finding reasons not to.

\sceneBreak

He went down to it alone, in the grey light, with the harbor already stirring above him into whatever confusion twenty enemy longships arriving at dawn was going to make of a city that had spent a year learning to expect only the dead from the north and nothing at all from the south.

Dragonbinder lay where it always lay, wrapped and chained, and Euron unwrapped it with hands that he noted, distantly, were not entirely steady — not from fear, he told himself, though he had long since stopped trusting the things he told himself in this particular cabin, in this particular dark. The horn was warm to the touch. It had not been warm before. He turned it in his hands and watched the valyrian glyphs along its length catch what little light the cabin permitted, and thought they moved rather more than poor light alone accounted for, sliding and reforming at the edge of his vision the way they always had, as though the thing were reading him back while he read it.

The tales were clear enough on the cost. He had made a point, back when he’d first taken the thing off dying men in the Grey Waste, of learning what the tales said, on the theory that a man ought to know the price of a thing before he committed to paying it. \textit{Dragonbinder,} the old Valyrian legends called it, \textit{hellhorn,} and every version of the tale agreed on the one particular the men who’d sold it to him had been too far gone to warn him of properly: that the horn bound dragons to the will of whoever sounded it, and that whoever sounded it — truly, with his whole chest behind the breath, not the soft testing he’d permitted himself once and paid three days’ agony for — burned for it. Not always at once. But always. The fire it commanded outside a man, the legends held, it first demanded inside him, and the binding held only so long as the one who’d sounded it lived to hold it, which in every telling was never very long at all.

He had spent his whole life refusing to pay other men’s prices. It was, he thought, turning the warm horn in his unsteady hands, rather a poor joke that the one price he found himself finally willing to pay was the steepest anyone had ever named him.

But Yara was in his harbor. The dead were five days out. The Golden Company would break, and Cersei would sit behind her walls until her walls came down, and every road out of this that Euron could see led through exactly one thing he could not out-patience or outspend or outlast — and a man who meant to still be standing at the end of the world, he understood at last, standing in the cold dark with the warm horn heavy in his hands, did not get to choose a cheaper way through it merely because he’d have preferred one.

“I’ve been called a great many things,” he said, aloud, to the horn, to the dark, to whatever vast winged thing the legends promised would turn in its sleep to look for him the moment he called it properly. “I don’t intend to add \textit{coward} to the list at the very end of it.”

He carried it up onto \textit{Silence}’s deck himself, into the grey dawn and the rising confusion, and stood a moment at her silent prow looking out at his niece’s sails and past them, north, toward a city that did not yet understand what he was about to do to the shape of its war. Then Euron Greyjoy filled his chest with the cold morning air, set the hellhorn to his lips, and sounded it with everything he had.

The sound was not a sound. That was the first thing, the only thing, that he had time to understand properly before the rest of it took him — that what came out of the horn was not heard so much as suffered, a note that went through a man rather than to him, and that somewhere beneath his own ribs something answered it that had no business being able to answer anything, a heat kindling in his chest exactly where the legends had promised it would kindle. He held the note. He held it past the point where holding it stopped being a thing he chose and became a thing the horn was doing to him, past the point where he tasted copper and then something worse than copper at the back of his throat, past the point where the edges of the grey harbor began to darken and narrow toward a single bright point directly ahead of him where, very far off, across the whole breadth of the doomed and freezing continent, something enormous turned slowly in the dark and, for the first time, looked back.

He did not stop until the horn stopped for him.

When it was done, Euron Greyjoy stood swaying at the prow of his silent ship with blood on his lips and a fire lit somewhere inside him that he understood, with the terrible clarity of a man who has finally paid a price he spent his whole life refusing, was never going to go out again for however many days he had left to carry it. His crew watched him, tongueless, unable to ask. He wiped the blood from his mouth with the back of one unsteady hand and found that he was, despite everything, despite the fire and the copper taste and the certain knowledge of what he’d just spent, smiling.

“There,” he said, to no one, to all of them, to the vast turning thing far off in the north that had finally, finally looked back. “Let’s see which of them can out-patience \textit{that.}”